\section{问题分析}

\subsection{总体分析}

四个子问题围绕“覆盖宽度”与“测线布设”两个主题层层递进：问题一建立二维剖面的基础几何模型，问题二将其推广到三维坡面，问题三在规则坡面海域上做最短测线优化，问题四进一步面向真实水深网格做数据驱动的分带布线。各问的覆盖宽度模型保持一致口径，即水平投影覆盖宽度，从而可与水平测线间距直接比较。

\subsection{问题一的分析}

问题一的核心是由换能器开角、坡度和测线位置确定单个条带的覆盖宽度及相邻条带的重叠率。由于题目给定的是理想斜坡剖面，所有参数已知，本问本质是二维几何建模与确定性计算。测线两侧的声束与坡面的交点不对称，左右半覆盖宽度不相等，因此重叠率不能直接套用平底公式 $\eta=1-d/W$，而应先计算相邻条带的代数重叠宽度，再除以当前条带覆盖宽度。

\subsection{问题二的分析}

问题二把二维剖面扩展到三维矩形海域。测线方向改变会影响两个量：船位处的水深，以及垂直测线剖面内的等效坡度。把坡面法向在水平面的投影方向与测线方向的夹角记为 $\beta$，可将三维问题通过向量投影转化为二维剖面问题，复用问题一的覆盖宽度公式。

\subsection{问题三的分析}

问题三的待求量是一组测线位置，目标是总长度最短，约束是完全覆盖和重叠率在 $10\%\sim20\%$。在单一平面坡面上，等深线为南北向直线，沿等深线布设时每条测线上的水深恒定，覆盖宽度仅随东西坐标变化，问题可降为一维布线；每条线贯穿矩形后长度固定，总长度最小等价于测线数最少。因此采用贪心递推，每步取最大可行间距，即令相邻重叠率取下限 $10\%$。

\subsection{问题四的分析}

问题四的水深不再是理想平面，而是规则网格数据。由于主方案采用南北向测线，覆盖宽度主要受东西向局部坡度影响，故由附件水深估计东西向梯度，建立局部线性海底剖面模型。为兼顾浅水与深水差异，将海域按南北方向分带，逐带递推布设线段，并在原始网格上仿真覆盖与重叠，比较不同分带粒度后选择可执行性较好的方案。
